\section{Multi-domain Task Incremental Learning}
\label{sec:setting}

\begin{figure}[t]
    \centering
    \includegraphics[width=\columnwidth]{figs/benchmark.pdf}
    \caption{Fig.\textbf{(a)}: examples of tasks from different domains in MTIL benchmark. Fig.\textbf{(b)}: illustration of calculating metrics Transfer, Avg. and Last during continual learning.}
    \label{fig:benchmark}
\end{figure}

\begin{table*}[t]
\centering
\caption{
\textbf{Ablation experiments}. 
% If not specified, the default setting is: the continual learning loss is image-only distillation, the teacher model is the initial CLIP model, the text source is ImageNet classes with a simple template, and the reference images are 100k sampled from ImageNet. 
Default settings are marked in \colorbox{baselinecolor}{gray}, which uses image and text distillation loss with the initial CLIP model on 100k ImageNet images and texts generated from ImageNet classes with a simple template.
}
%#################################################
% loss type
%#################################################
\subfloat[
    \textbf{Continual learning loss}.
    \label{tab:loss_type}
]{
    \centering
    \begin{minipage}{0.29\linewidth}{
        \begin{center}
            \tablestyle{1pt}{1.05}
            \begin{tabular}{y{50}x{30}x{30}x{30}}
                loss & Transfer & Avg. & Last \\
                \shline
                CE only & 44.6 & 55.9 & 77.3 \\
                Feat. Dist. & 47.6 & 58.7 & 77.1 \\
                Image-only & 56.5 & 68.9 & 82.1 \\
                Text-only & 56.7 & 69.0 & 82.6 \\
                Both & \baseline{\textbf{56.8}} & \baseline{\textbf{69.2}} & \baseline{\textbf{83.0}}
        \end{tabular}
        \end{center}
    }\end{minipage}
}
\hspace{2em}
%#################################################
% Data Source
%#################################################
\subfloat[
\textbf{Data sources for replay}.
\label{tab:image_source}
]{
\centering
\begin{minipage}{0.29\linewidth}{\begin{center}
\tablestyle{4pt}{1.05}
\begin{tabular}{y{30}x{30}x{30}x{30}}
source & Transfer & Avg. & Last \\
\shline
    current & 56.7 & 66.5 & 80.2 \\
    ImageNet & \baseline{56.8} & \baseline{\textbf{69.2}} & \baseline{\textbf{83.0}} \\
    CC & \textbf{57.2} & 68.5 & 80.9 \\
    CIFAR100 & 55.2 & 65.9 & 80.7 \\
    Flowers & 54.7 & 66.0 & 80.8 \\
\end{tabular}
\end{center}}\end{minipage}
}
\hspace{2em}
%#################################################
% Text Source
%#################################################
\subfloat[
    \textbf{Text sources for replay}.
    \label{tab:text_source}
]{
    \centering
    \begin{minipage}{0.29\linewidth}{
        \begin{center}
            \tablestyle{1pt}{1.05}
            \begin{tabular}{y{54}x{30}x{30}x{30}}
                source & Transfer & Avg. & Last \\
                \shline
                current & 51.8 & 64.9 & 82.0 \\
                prev. all & 54.0 & 70.2 & 83.7 \\
                1k classes (IN) & \baseline{56.8} & \baseline{69.2} & \baseline{83.0} \\
                13k Sent. (CC) & \textbf{58.9} & \textbf{70.5} & \textbf{84.0} \\
                1k Rand. Sent. & 58.7 & 70.2 & 83.8 \\
                % 100 Random & 61.8 & \textbf{90.1} & 86.7 \\
            \end{tabular}
        \end{center}}\end{minipage}
}
%#################################################
\\
\centering
%#################################################
% Teacher model
%#################################################
\subfloat[
\textbf{Teacher model}.
\label{tab:ref_model}
]{
\centering
\begin{minipage}{0.29\linewidth}{\begin{center}
\tablestyle{4pt}{1.05}
\begin{tabular}{x{30}x{30}x{30}x{30}}
source & Transfer & Avg. & Last \\
\shline
    Initial & \baseline{\textbf{56.8}} & \baseline{\textbf{69.2}} & \baseline{\textbf{83.0}} \\
    $n-1$ & 53.9 & 66.6 & 80.7 \\
    WiSE(0.5) & 56.4 & 68.9 & 82.9 \\
    WiSE(0.8) & 56.2 & 67.8 & 81.3 
\end{tabular}
\end{center}}\end{minipage}
}
\hspace{2em}
%#################################################
% Replay image
%#################################################
\subfloat[
\textbf{\# images for replay}.
\label{tab:replay_num}
]{
\begin{minipage}{0.29\linewidth}{\begin{center}
\tablestyle{4pt}{1.05}
\begin{tabular}{x{30}x{30}x{30}x{30}}
\#image & Transfer & Avg. & Last \\
\shline
% 1.28M &62.9 &89.7 &86.3 \\
1M & \textbf{58.7} & \textbf{70.1} & \textbf{83.2} \\
100k & \baseline{56.8} & \baseline{69.2} & \baseline{83.0} \\
10k & 57.8 & 68.7 & 81.2 \\
1k & 56.3 & 67.6 & 80.8 \\
\end{tabular}
\end{center}}\end{minipage}
}
\hspace{2em}
%#################################################
% Replay classes
%#################################################
\subfloat[
\textbf{\# image classes for replay}.
\label{tab:replay_classes}
]{
\begin{minipage}{0.29\linewidth}{\begin{center}
\tablestyle{4pt}{1.05}
\begin{tabular}{x{30}x{30}x{30}x{30}}
\#class & Transfer & Avg. & Last \\
\shline
1000 & \baseline{\textbf{56.8}} & \baseline{\textbf{67.6}} & \baseline{\textbf{83.0}} \\
100 & 56.7 & 67.3 & 82.3 \\
10 & 53.8 & 66.4 & 81.0 \\
1 & 53.1 & 65.5 & 80.5 
\end{tabular}
\end{center}}\end{minipage}
}
\label{tab:ablations}
\end{table*}

\paragraph{Conventional Continual Learning Benchmark.} A benchmark consisting of several tasks is needed to evaluate different methods for continual learning. Most previous benchmarks are built by separating classes in a single dataset, such as MNIST~\cite{van2019three}, CIFAR100~\cite{douillard2022dytox}, TinyImageNet~\cite{yan2021dynamically}, and ImageNet~\cite{yan2021dynamically,zhu2021prototype}. We also evaluate our method with traditional benchmarks. In CIFAR100~\cite{Krizhevsky09learningmultiple}, classes are separated into groups to build tasks. Suppose the dataset has $m$ classes, a $k$-step setting means we learn $\nicefrac{m}{k}$ classes in each new task. The CIFAR100 contains 100 classes, and the setting of 10, 20, and 50 steps are visited. For TinyImageNet with $m=200$, the first step learns 100 classes, and the rest is learned with 5, 10, and 20 steps. As for ImageNet-100, we have two settings: ImageNet-100-B0, which includes the same amount of classes for each step, and ImageNet-100-B50, which has 50 classes for the first step, and the remaining 50 classes are observed progressively over the next 10 stages. For ImageNet~\cite{deng2009imagenet}, we investigate a 10-step setting, which learns 100 new classes per task.

\paragraph{MTIL Benchmark.} Different classes from one dataset share the common image source and a similar style~\cite{recht2019imagenet,hendrycks2021many}. Thus, we propose Multi-domain Task Incremental Learning (MTIL), a cross-domain version of task incremental learning. Different tasks are collected from different domains, requiring different domain knowledge for humans to achieve high accuracy. Our MTIL benchmark consists of 11 tasks (detailed in supplementary materials), as some of the tasks illustrated in \cref{fig:benchmark} (a). The MTIL benchmark is very challenging with a total number of 1,201 classes. Two orders are used for the evaluation; the first one is alphabet order (Order-I): Aircraft, Caltech101, CIFAR100, DTD, EuroSAT, Flowers, Food, MNIST, OxfordPet, StanfordCars, SUN397. And the second one is a random order (Order-II): StanfordCars, Food, MNIST, OxfordPet, Flowers, SUN397, Aircraft, Caltech101, DTD, EuroSAT, CIFAR100. Experiments are done in Order I by default.


\begin{table}[t]
\centering
\caption{Ablation study of different components for ZSCL.}
\label{tab:pablate}
\resizebox{\columnwidth}{!}{%
% {\small
\begin{tabular}{@{}lcc|cc|cc@{}}
\toprule
Method & Transfer & $\Delta$ & Avg. & $\Delta$ & Last & $\Delta$  \\ \midrule
CLIP \texttt{ViT-B/16@224px}  & & & & & & \\
\quad Zero-shot  & \baseline{69.4} &  \baseline{0.0} & \baseline{65.3} & \baseline{0.0} & \baseline{65.3} & \baseline{0.0} \\
\quad Continual Learning & 44.6 & \red{-24.8} & 55.9 & \red{-9.4} & 77.3 & \gre{+12.0} \\
\quad + Distillation & 58.9 & \red{-10.5} & 70.5 & \gre{+5.2} & 83.8 & \gre{+18.5}  \\
\quad \quad + WiSE-FT (best $\alpha$) & 61.7 & \red{-7.7} & 71.6 & \gre{+6.3} & 83.3 & \gre{+18.0}\\
% \quad \quad + WE (ZSCL) & 65.8 & \red{-1.6} & \textbf{89.7} & \gre{+0.7 }& \textbf{86.8} & \gre{+19.4} \\
\quad \quad + WE (ZSCL$^*$) & 62.2 & \red{-7.2} & 72.6 & \gre{+7.3} & \textbf{84.5} & \gre{\textbf{+19.2}} \\
\quad \quad + WC & 67.6 & \red{-1.8} & 74.5 & \gre{+9.2} & 83.2 & \gre{+17.9}\\
\quad \quad \quad + WiSE-FT & 67.7 & \red{-1.7} & 74.2 & \gre{+8.9} & 81.9 & \gre{+16.6} \\
\quad \quad \quad + WE (ZSCL) & \textbf{68.1} & \red{\textbf{-1.3}} & \textbf{75.4} & \gre{\textbf{+10.1}} & 83.6 & \gre{+18.3} \\
\bottomrule
\end{tabular}%
}
\end{table}


\paragraph{Evaluation Metrics.} The metrics of MTIL are illustrated in \cref{fig:benchmark}(b), where rows represent training steps and a column shows performance for one dataset.
For conventional continual learning, only scores under the diagonal are meaningful, since they cannot give zero-shot predictions on unseen tasks. In contrast, the zeros-hot transfer ability enables a vision-language model to provide predictions for all datasets. The average accuracy on all datasets among all timestamps is the ``Avg'' metric. The ``Last'' metric is the average performance of all tasks after CL. The ``Transfer'' metric is the average task performance in the upper-right triangle of the matrix. Every task's performance is first averaged to equal the weight of each dataset. It measures to what extent the zero-shot transfer ability is preserved. Before learning task $i$, tasks not earlier than $i$ are not fine-tuned. Thus, their performance is an indicator of zero-shot transfer ability. 